\documentclass[11pt]{article}
\usepackage[a4paper,margin=20mm]{geometry}
\usepackage{amsmath,amssymb,amsthm,microtype}
\usepackage[colorlinks=true,linkcolor=blue,urlcolor=blue]{hyperref}
\newtheorem{lemma}{Lemma}
\newtheorem{theorem}[lemma]{Theorem}
\newtheorem{corollary}[lemma]{Corollary}
\title{Exact rank-two diagonal MPU circuits and approximation at arbitrary rank}
\author{TNLean contributors}
\date{2 October 2026}
\begin{document}
\maketitle

Let $f:\{0,1\}^N\to\mathbb C$ satisfy $|f(x)|=1$, and let
$U_f|x\rangle=f(x)|x\rangle$. At cut $k$, write
$M_k(x,y)=f(xy)$, where $x$ has length $k$. The operator Schmidt rank of
$U_f$ at this cut is exactly $\operatorname{rank}M_k$. Thus an
open-boundary diagonal MPO of maximal bond dimension $D$ gives
$\operatorname{rank}M_k\leq D$. All circuits below use arbitrary one- and
two-qubit gates; auxiliary qubits begin and end in $|0\rangle$.

The general MPU question in Styliaris--Trivedi--Cirac,
arXiv:2508.08160, concerns arbitrary unitaries, including those which
mix computational basis vectors. The results here concern diagonal unitaries,
and, by a separate reduction, monomial unitaries. They remove the
smallest-Schmidt-value hypothesis within these classes. The general fixed-bond
exact polynomial circuit theorem, and its operator-norm approximation
corollaries with exact auxiliary erasure, are proved in the companion note
\emph{Exact polynomial circuits for matrix product unitaries}, Section~5
(\nolinkurl{docs/audits/2026-10-02_mpu_rank_two_circuits.tex}).
No claim of priority is intended. The relevant source
passages of arXiv:2508.08160v2 are the introductory discussion of
circuit implementations of MPUs and the main-text paragraph on nonuniform
MPUs containing the definition of $q_k$ and Theorem~2 (lines 683--690
and 1373--1402 of \texttt{main.tex} in the version-2 arXiv source,
\url{https://arxiv.org/e-print/2508.08160v2}).

\section{The rank-two dichotomy}

\begin{lemma}\label{lem:circle}
Suppose $M$ is a finite, nonempty matrix, all its entries have modulus one,
and $\operatorname{rank}M\leq2$. Either its rows have at most two types
up to multiplication by a nonzero scalar, or its columns have at most two
such types.
\end{lemma}
\begin{proof}
Rank one gives a single row type. Otherwise choose two independent rows
$u,v$. Divide each column by its corresponding entry of $u$. Every row
then has the form $a_x+b_x t_y$, with $t_y=v_y/u_y$ of modulus one.
If $t$ takes at most two values, there are at most two column types.
If it takes at least three values, the identity
\[
 1=|a_x+b_x t_y|^2
   =|a_x|^2+|b_x|^2+2\operatorname{Re}(a_x\overline{b_x}\overline{t_y})
\]
forces $a_x\overline{b_x}=0$: a nonconstant real linear functional
on the plane has a level line, and that line meets the unit circle in
at most two points. Thus each row is a scalar multiple of either $1$ or
$t$, giving at most two row types.
\end{proof}

The circle--line argument, its matrix-basis step, and the circuit
construction below are mathematical proofs; their Lean formalizations
are not claimed in this note.

\begin{lemma}\label{lem:selector}
Suppose all consecutive cut ranks of a unit-modulus function $f(xy)$ are
at most two, and there are two projective row types across the displayed
cut. There are phase functions $g(x),r_0(y),r_1(y)$ and a Boolean function
$p(x)$ such that
\[
 f(xy)=g(x)r_{p(x)}(y).
\]
Each of $g,r_0,r_1$ has consecutive cut ranks at most two.
The Boolean function $p$ has consecutive cut ranks at most five, and
can be computed reversibly with $O(|x|)$ gates and auxiliary qubits.
It also admits a reversible computation of depth $O(\log(|x|+1))$
using $O(|x|)$ gates and auxiliary qubits.
\end{lemma}
\begin{proof}
Fix $y_0$. Normalize representative rows so that $r_s(y_0)=1$ and put
$g(x)=f(xy_0)$. Then $r_s(y)=f(x_s y)/f(x_s y_0)$ for suitable
representatives $x_s$. These three functions are restrictions of $f$,
up to a constant phase, so their internal cut ranks remain at most two.
Choose $y_1$ with $r_0(y_1)=\alpha\ne\beta=r_1(y_1)$. Then
\[
 p(x)=\frac{f(xy_1)\overline{f(xy_0)}-\alpha}{\beta-\alpha}.
\]
At an internal cut, the product in the numerator has rank at most
$2\cdot2=4$; subtraction of a constant raises rank by at most one.
Thus $p$ has cut ranks at most five, regardless of how small
$|\beta-\alpha|$ is.

A zero-one coefficient table of rank $r$ has at most $2^r$ distinct rows:
choose $r$ columns spanning the column space; those columns distinguish
all rows. Consequently $p$ has at most 32 residual functions at every
cut. Reading a bit changes the residual function by a deterministic
transition. The resulting Boolean computation uses a fixed finite set
of states. Each transition is a function on at most 32 states, selected
by the input bit. Composition of these functions is associative and
can be computed by a balanced tree with $O(|x|)$ operations and depth
$O(\log(|x|+1))$. Each operation has constant size. Its Boolean circuit
can be made reversible by retaining intermediate results, copying the
final answer, and reversing the computation. The cost remains linear,
and the depth remains logarithmic.
\end{proof}

The finite-alphabet row bound is proved in
\nolinkurl{TNLean/MPS/MPU/FiniteAlphabetRows.lean}. The affine-product
factorization and Boolean residual bound are established in the preceding
mathematical proof. The denominator in the selector formula identifies a
Boolean function classically; no quantum amplitude amplification is used.

\section{An exact circuit}

\begin{theorem}\label{thm:exact}
If every consecutive operator Schmidt rank of a diagonal $N$-qubit
unitary is at most two, it has an exact circuit with
\[
 G=O(N^{\log_2 3})
\]
one- and two-qubit gates. One may choose either $O(N)$ auxiliary qubits
and depth $O(G)$, or $O(G)$ auxiliary qubits and depth $O(\log(N+1))$.
On a line, the first construction can be routed to give gate count and
depth $O(N^{1+\log_2 3})$, with $O(N)$ auxiliary qubits.
\end{theorem}
\begin{proof}
Pad to a power of two by adding unused variables; this changes $N$ by a
factor less than two. Split the variables into equal consecutive halves.
Lemma~\ref{lem:circle} and its transpose give a decomposition as in
Lemma~\ref{lem:selector}, possibly with the two halves interchanged.
A rank-one cut requires only two phase factors.

It is convenient to implement a phase function controlled by a single
bit $c$. Compute $p$ reversibly, and form $c_0=c\wedge\neg p$ and
$c_1=c\wedge p$. Recursively apply $g$ controlled by $c$, $r_0$
controlled by $c_0$, and $r_1$ controlled by $c_1$. Erase the two new
controls and $p$. On every computational basis vector the resulting
phase is $g(x)^c r_{p(x)}(y)^c$, as required. By linearity the same
identity holds on an arbitrary input state. At a leaf, an arbitrary
single-variable phase under one control is a two-qubit diagonal gate.
The root control is initialized to one and erased at the end.

The gate count satisfies $G(n)\leq3G(n/2)+Cn$, so
$G(n)=O(n^{\log_2 3})$. In a sequential traversal, auxiliary storage is
reused between the three recursive calls. The simultaneously needed
interval lengths sum to at most $n+n/2+n/4+\cdots$, giving $O(n)$
auxiliary qubits.

For the depth bound, give every node its own copies of the computational
basis bits in its interval, using controlled-NOT gates. This operation
is coherent on superpositions; it does not copy an arbitrary unknown
quantum state. The total number of copies is bounded by
\[
 \sum_{d=0}^{\log_2N}3^d\frac{N}{2^d}=O(N^{\log_2 3}).
\]
All copies can be distributed in depth $O(\log N)$ with arbitrary
pairwise interactions. Compute all node selectors in parallel by the
balanced computations in Lemma~\ref{lem:selector}. Propagate the controls
from root to leaves, in constant depth per tree level. Apply the leaf
phase gates in parallel, then reverse the controls, selectors, and
copies. The total depth is logarithmic and all auxiliary qubits are
erased. Finally, in the sequential construction, at most $O(N)$ qubits
are present. Routing each two-qubit interaction by $O(N)$ swaps gives
the stated bound on a line.
\end{proof}

The proof is an existence statement for arbitrary exact gates. It does
not bound the bit complexity of identifying the two projective types
from inexact tensor data. Nor does it give exact synthesis over a fixed
finite gate set. These distinctions are essential when phases can be
arbitrary complex numbers.

\section{Approximation without a rank-two restriction}

\begin{theorem}\label{thm:approx}
Suppose $\operatorname{rank}M_k\leq D$ at every consecutive cut.
For $0<\varepsilon<1$, the diagonal unitary $U_f$ has a circuit whose
restriction to data with clean auxiliary input is
$\widetilde U_f\otimes|0\rangle$, where
\[
 \|\widetilde U_f-U_f\|\leq\varepsilon.
\]
Its gate count and auxiliary space are polynomial in
$N,D,\log(1/\varepsilon)$. No lower bound on a nonzero Schmidt value is
needed. This is a nonuniform existence statement; a precision or
conditioning guarantee for converting arbitrary supplied tensor data
is not asserted.
\end{theorem}
\begin{proof}
Normalize the phase vector as
$\psi(x)=2^{-N/2}f(x)$. Its norm is one and its consecutive Schmidt
ranks are at most $D$. Choose a right-canonical open-boundary MPS for
$\psi$. Write its matrices as $A_j(x_j)$; each has operator norm at most
one, since the canonical condition is
$\sum_s A_j(s)A_j(s)^\dagger=I$. The two endpoint bonds have dimension
one. Hence
\[
 \psi(x)=A_1(x_1)\cdots A_N(x_N),\qquad \|A_j(s)\|\leq1.
\]
Round every real and imaginary matrix entry to $b$ binary fractional
places, and evaluate the product with fixed-point arithmetic, rounding
each intermediate scalar operation to the same precision. There are
$O(ND^2)$ such operations in a row-vector contraction. A matrix-entry
error gives operator-norm error at most $2D2^{-b}$. At each step the
rounding of the vector contraction has Euclidean error at most
$CD^2 2^{-b}$ for an absolute $C$. Comparing with the exact product,
and using the norm-one bound on the exact matrices, gives
\[
 |\widehat\psi(x)-\psi(x)|\leq C'ND^2 2^{-b}
\]
uniformly in $x$, provided $ND^2 2^{-b}$ is sufficiently small.
For example this follows by iterating
$e_{j+1}\leq(1+2D2^{-b})e_j+CD^2 2^{-b}$.
Intermediate vectors stay bounded, so $O(\log(D+1))$ integer guard bits
suffice, including the unrounded sums.

Take
\[
 b=\left\lceil N/2+\log_2(C''ND^2/\varepsilon)\right\rceil
\]
with a sufficiently large absolute constant $C''$. Rescaling by
$2^{N/2}$ then computes a number $v(x)$ with
$|v(x)-f(x)|\leq\varepsilon/20$. For odd $N$, the additional factor
$\sqrt2$ can be rounded to the same accuracy. Thus $|v(x)|\geq19/20$.
Computing its argument to error at most $\varepsilon/10$ requires
only polynomially many bit operations. One explicit procedure chooses
the larger absolute component, divides the smaller by it, and reduces
the argument to $\arctan t$ with $|t|\leq1$. The identity
\[
 \arctan t=2\arctan\frac{t}{1+\sqrt{1+t^2}}
\]
reduces the Taylor-series argument in absolute value to at most
$\sqrt2-1$. The series therefore converges geometrically. Division has
a denominator bounded away from zero; the square root is in $[1,\sqrt2]$.
Binary division, bisection for the square root, and
$O(\log(1/\varepsilon))$ Taylor terms, with suitably increased working
precision, all have polynomial cost. Signs and the chosen component
determine the quadrant. A possible discontinuity of the principal
argument across its branch cut is immaterial after exponentiation.

Implement this Boolean calculation reversibly, retaining its intermediate
results. If its final angle is represented in binary, apply the
corresponding single-qubit phase gates to those bits. Their product
multiplies the basis input by the computed phase. Reverse the entire
calculation. Auxiliary qubits return exactly to zero. The resulting
phase $\widetilde f(x)$ obeys
$|\widetilde f(x)-f(x)|\leq\varepsilon$ uniformly: normalization of
$v(x)$ contributes at most $2|v(x)-f(x)|$, and the angle approximation
contributes its absolute error. For diagonal operators, the operator
norm of their difference is precisely this maximum phase error.
All working precisions and operation counts are polynomial in the
stated parameters, completing the proof.
\end{proof}

The factor $2^{N/2}$ costs $O(N)$ precision bits, rather than an
exponential number of gates. This is the distinction between evaluating
a diagonal phase and implementing a general unitary from its matrix
entries. The latter cannot be reduced to phase evaluation in this way.

\begin{corollary}\label{cor:monomial}
Every monomial MPU of fixed maximal bond dimension admits an
operator-norm approximation with gate count and clean auxiliary space
polynomial in $N,\log(1/\varepsilon)$, without restricting its phase
alphabet or smallest Schmidt value.
\end{corollary}
\begin{proof}
Write $U|x\rangle=f(x)|F(x)\rangle$. The permutation matrix $P_F$ has
coefficients $|U_{yx}|^2$, so it has an MPO representation of maximal
bond dimension at most $D^2$. The diagonal unitary $P_F^\dagger U$
then has maximal bond dimension at most $D^3$. The exact permutation
circuit proved in
\nolinkurl{docs/audits/2026-10-02_mpu_finite_alphabet_circuits.tex}
implements $P_F$ with $O(N8^{D^2})$ gates. Apply
Theorem~\ref{thm:approx} to the diagonal factor. Left multiplication
by $P_F$ preserves operator-norm error. The dependence on variable $D$
in this corollary is exponential, because of the permutation step.
\end{proof}

\section{Verification and relation to the general theorem}

The script \nolinkurl{scripts/check_mpu_diagonal_rank_two.py} checks
exact recursive reconstruction on nearest-neighbour phase interactions,
star interactions in both orientations, parity phases, projection phases, and product phases,
using rational complex numbers of modulus one. It verifies all cut ranks,
all Boolean selector rank bounds, and the residual-state bound on each
recursive call. Enumeration of small examples is a check of the
identities, not a proof of the asymptotic circuit bound.

The approximation theorem here is a mathematical argument, not a completed
Lean proof or an implemented reversible arithmetic circuit. The exact
circuit theorem here is likewise a separate mathematical statement from
the formalized finite-alphabet row bound. For arbitrary MPUs of fixed bond dimension, the
companion note's Section~5 proves an exact polynomial circuit bound and
its operator-norm approximation corollaries, with exactly erased auxiliary
registers. A Lean formalization of those general statements is pending in
TNLean pull request~\#8590. This does not
assert a formalization of the sharper diagonal circuit bounds or of the
reversible-arithmetic construction proved in this note.
\end{document}
